\chapter{Experimental Design}

\section{Planning and Conducting Experiments}

Surveys tell you what people think or report. Experiments tell you what actually causes what. 
The distinction  determines what conclusions you are entitled to draw when the 
data comes in.

\subsection{Experiments vs. Observational Studies}

The defining question is whether the researcher assigns the conditions being studied or simply observes them as they exist.

In an \newterm{observational study}\index{observational study}, the researcher collects data on subjects without intervening. 
The variables of interest already exist in the subjects; the researcher records what is there. A sociologist who tracks the grades 
and sleep hours of university students over a semester, without changing anyone's sleep schedule, is conducting an observational study.

In an \newterm{experiment}\index{experiment}, the researcher imposes conditions on subjects and then measures the 
effect or change. The researcher creates the differences rather than observing them. 

Why does the distinction matter? Only a properly designed experiment can support a cause-and-effect 
conclusion. An observational 
study can show that two variables are associated, but a third variable you did not account for could 
explain the relationship 
entirely. This uninvited third variable is called a 
\newterm{confounding variable}\index{confounding variable}: 
one whose effect on the 
response cannot be separated from the effect of the variable you are actually studying.

% FIXME image/illustration needed

In the observational sleep study, students who sleep more might also exercise more, eat better, or have lower course loads. 
Any of these could be responsible for better grades. There is no way to figure out which effects correlate to 
which behaviours. 

There are situations where an experiment is unethical or impossible: you cannot randomly assign people to smoke for twenty years, 
or randomly select cities to experience floods. In these cases, observational data is all we have, and conclusions must be drawn 
with caution.

\begin{Exercise}[title={Experiment or Observational Study?}, label={ex:observe}]
For each study, identify whether it is an experiment or an observational study. Then state whether a cause-and-effect 
conclusion is justified and why.

\medskip
\begin{tabular}{|p{0.53\textwidth}|p{0.17\textwidth}|p{0.19\textwidth}|}
\hline
\textbf{Study} & \textbf{Type} & \textbf{Causal conclusion?} \\
\hline
Researchers randomly assign 180 seedlings to three groups and water each group with a different fertilizer solution. 
After eight weeks they measure plant height. & & \\[2.5em]
\hline
A researcher records the screen time and attention span of 500 children aged 8 to 12, then looks for a relationship 
between the two. & & \\[2.5em]
\hline
A company randomly selects half its customer service team to complete a new communication training program, 
then compares customer satisfaction scores between the two groups. & & \\[2.5em]
\hline
A historian finds that cities with more libraries per capita had higher literacy rates in the early 20th century. & & \\[2.5em]
\hline
\end{tabular}
\end{Exercise}
\begin{Answer}[ref={ex:observe}]
\begin{itemize}
    \item \textbf{Seedlings:} Experiment. Fertilizer was randomly assigned. Causal conclusion is justified.
    \item \textbf{Screen time and attention:} Observational study. No assignment was made. Causal conclusion is not 
    justified; many confounding variables (parenting style, diet, sleep) could explain the relationship.
    \item \textbf{Customer service training:} Experiment. The training was randomly assigned. Causal conclusion is justified.
    \item \textbf{Libraries and literacy:} Observational study. No assignment was made. Causal conclusion is not justified; 
    wealthier cities may have had both more libraries and higher literacy for unrelated reasons.
\end{itemize}
\end{Answer}

\subsection{Vocabulary for Experiments}

Experiments have their own vocabulary. These terms build on each other, so it helps to learn them in order.

% FIXME image/illustration needed

\begin{description}
    \item[\newterm{Explanatory variable}\index{explanatory variable}] Also called the independent variable. The variable 
    whose effect on the response is being studied. In a study of whether LoFi music affects study focus, the type of content 
    being studied is the explanatory variable.

    \item[\newterm{Response variable}\index{response variable}] Also called the dependent variable. The outcome being measured. 
    In a study of whether background music affects study typing speed, the typing speed is the response variable.

    \item[\newterm{Experimental unit}\index{experimental unit}] The smallest unit to which a treatment is applied. In a study of 
    plants, each plant is an experimental unit. In a study of people, each person is an experimental unit.

    \item[\newterm{Factor}\index{factor}] A variable whose effect on the response is of interest. Factors can be qualitative 
    (type of music: classical, jazz, Afrobeat) or quantitative (temperature in degrees Kelvin).

    \item[\newterm{Levels}\index{levels}] The specific values of a factor used in the experiment. If the factor is type of music 
    with three options, there are three levels.

    \item[\newterm{Treatment}\index{treatment}] A specific combination of factor levels applied to experimental units. With one 
    factor and three levels, there are three treatments. With two factors at two levels each, there are $2 \times 2 = 4$ treatments.

    \item[\newterm{Control group}\index{control group}] A group that receives no treatment or a standard baseline condition. 
    It provides a reference point for comparison with the treatment groups.

    \item[\newterm{Placebo}\index{placebo}] A dummy treatment that resembles the real treatment but contains no active component. 
    Used to isolate the genuine effect of the treatment from the \newterm{placebo effect}\index{placebo effect}: the tendency for 
    people to show change simply because they believe they are receiving treatment.
\end{description}

\subsection{The Three Principles of Experimental Design}

A well-designed experiment rests on three principles. 

% FIXME image/illustration needed

\newterm{Control}\index{control!in experiments} means keeping everything the same for all groups except the treatment. It is not just having a control group. Control helps make sure the treatment is the only real difference.

\newterm{Randomization}\index{randomization} means using chance to place experimental units into groups. It helps keep the groups fair and spreads unknown factors across all groups.

A \newterm{lurking variable}\index{lurking variable} is something that can affect the response but is not part of the experiment.

\newterm{Replication}\index{replication} means including enough experimental units in each group so that chance variation 
among individuals does not swamp the treatment effect. If a new training method is tested on only two employees per group, 
a single unusually high or low performer could dominate the group's result. With 30 employees per group, individual extremes 
average out.

\begin{Exercise}[title={Applying the Three Principles}, label={ex:three-principles}]
A horticulture researcher wants to test whether exposing seedlings to classical music for two hours per day increases their 
growth rate compared to no music exposure. She has 60 seedlings of the same species, grown from the same seed batch, available for 
the study.

\begin{enumerate} 
    \item Describe how randomization should be applied.
    \item What does control mean in this context? List at least three variables that should be held constant.
    \item The researcher initially plans to use 5 seedlings per group. A colleague recommends 30 per group. Which principle does 
    this change address, and why?
    \item Should there be a control group? What would it look like?
\end{enumerate}
\end{Exercise}
\begin{Answer}[ref={ex:three-principles}]
\begin{enumerate} 
    \item Number the 60 seedlings from 1 to 60. Use a random number generator to select 30 for the music group. The remaining 
    30 form the no-music group.
    \item Control means growing all seedlings under identical conditions except for the music. Variables to hold constant include: 
    soil type and quantity, pot size, amount and frequency of watering, light exposure, temperature, and humidity.
    \item Replication. With only 5 seedlings per group, a single unusually fast or slow grower could distort the group average. 
    With 30 per group, individual variation averages out and a real effect of the music is more detectable.
    \item Yes. The no-music group serves as the control group, establishing the baseline growth rate without the treatment. 
    Without it, there is nothing to compare the music group's growth against.
\end{enumerate}
\end{Answer}


\subsection{Blinding}

When subjects know which treatment they are receiving, their behavior or reported experience may 
change in ways that have nothing to do with the treatment itself. When the person measuring outcomes 
knows which group a subject belongs to, their assessments may be influenced a littke. Both of these are 
sources of bias that blinding is designed to prevent.

In a \newterm{single-blind experiment}\index{single-blind experiment}, either the subjects do not 
know which treatment they received, or the person measuring outcomes does not know, but not both. 

In a \newterm{double-blind experiment}\index{double-blind experiment}, neither the subjects nor the 
evaluators know the treatment assignments. Double-blind experiments are preferred when feasible because 
they eliminate bias from both sources simultaneously.

Double-blinding is not always possible. 
If subjects can easily identify their treatment condition from its appearance or 
feel, blinding them is impossible. A study comparing two different workout splits cannot prevent 
participants from knowing which one they are following. In such cases, the best available design is 
to blind the evaluators, even if the subjects cannot be blinded.

\subsection{Random Selection vs. Random Assignment}

What you are allowed to conclude from the results depends on two choices made at the design stage: 
whether subjects were randomly assigned to treatments, and whether subjects were randomly selected 
from a larger population. These are two separate decisions, and they lead to two separate types of conclusions.

\newterm{Random assignment}\index{random assignment} of treatments is what allows you to conclude that a
 treatment caused an observed difference. When treatments are assigned by chance, any pre-existing differences 
 among subjects are distributed randomly across groups. If one group then performs meaningfully better than 
 another, the only systematic explanation is the treatment itself. A difference large enough to be unlikely 
 due to chance alone is called \newterm{statistically significant}\index{statistically significant}. 
 Statistically significant differences between treatment groups are evidence that the treatment caused 
 the effect.

\newterm{Random selection}\index{random selection} of subjects from a population is what allows you to 
generalize conclusions beyond the people in your study. If subjects were randomly drawn from a well-defined 
population, the results of the experiment can be extended to that population. If subjects were not randomly 
selected, the conclusions apply only to the subjects in the study, or to others who are very similar to them.

These two ideas operate independently. An experiment can have one, both, or neither. The table below summarizes what each combination allows.

\begin{center}
\begin{tabular}{|p{0.22\textwidth}|p{0.22\textwidth}|p{0.22\textwidth}|p{0.22\textwidth}|}
\hline
 & \textbf{Random assignment: Yes} & \textbf{Random assignment: No} \\
\hline
\textbf{Random selection: Yes} & Can conclude causation. Can generalize to the population. & Cannot conclude causation. Can generalize to the population. \\
\hline
\textbf{Random selection: No} & Can conclude causation. Cannot generalize beyond the study subjects. & Cannot conclude causation. Cannot generalize beyond the study subjects. \\
\hline
\end{tabular}
\end{center}

Remember: These two are not interchangeable! However, you need both to make a causal claim that applies to a broader population.

\begin{Exercise}[title={Random Assignment vs. Random Selection}, label={ex:assign-vs-select}]
For each study, answer three questions: (1) Was there random assignment? (2) Was there random selection? (3) What conclusions are justified?

\begin{enumerate} 
    \item A researcher posts a sign-up sheet in a university cafeteria asking for volunteers to participate in a study on the effect of background noise on reading comprehension. Sixty students volunteer. The researcher randomly assigns 30 to read in a noisy room and 30 to read in a quiet room. Students in the noisy room score significantly lower on a comprehension test.

    \item A public health agency randomly selects 400 adults from the voter registration rolls of a large city. Researchers observe whether each person exercises regularly and whether they report high stress levels. Those who exercise regularly report significantly lower stress.

    \item A school district randomly selects 10 of its 60 elementary schools to pilot a new reading curriculum. All students in those 10 schools use the new curriculum; the remaining 50 schools continue with the standard curriculum. At year end, students in the pilot schools score significantly higher on a standardized reading test.

    \item Nike recruits 80 volunteers through social media and randomly assigns 40 to a new training app and 40 to a standard printed workout plan. After 12 weeks, the app group shows significantly greater improvement in endurance.
\end{enumerate}
\end{Exercise}
\begin{Answer}[ref={ex:assign-vs-select}]
\begin{enumerate} 
    \item (1) Yes, random assignment. (2) No, volunteers are self-selected. Conclusion: the noisy room likely caused lower comprehension scores, but this result cannot be generalized to all university students or beyond. Volunteers may differ from the broader student population in ways that affect how noise impacts their reading.

    \item (1) No random assignment because the researcher observed existing behavior. They did not observe an assigned treatment. (2) Yes, random selection from voter rolls. Conclusion: the association between exercise and lower stress can be generalized to the adult population of the city. We cannot conclude that exercise caused lower stress because no treatment was assigned. A confounding variable such as income or work conditions could explain both.

    \item (1) No random assignment of individual students to the curriculum because schools were the units assigned rather than students. The assignment was at the school level. (2) Yes, random selection of schools from the district. Conclusion: results can be generalized to the district's elementary schools. The random assignment of schools to curricula supports a causal claim at the school level, but student-level confounding (differences in which students attend which schools) is a concern worth noting.

    \item (1) Yes, random assignment to app or printed plan. (2) No, volunteers are recruited through social media. Conclusion: the app likely caused greater endurance improvement compared to the printed plan. However, this result cannot be generalized to all people interested in fitness, as social media recruits may be younger, more comfortable with mobile phones and social media in general, or more motivated than the broader population.
\end{enumerate}
\end{Answer}



\section{Experimental Designs}

There are three experimental designs you should be able to recognize, describe, and distinguish from one another. Each one handles variability among experimental units differently.

\subsection{Completely Randomized Design}

The simplest structure is a \newterm{completely randomized design}\index{completely randomized design}. All available experimental units are assigned to treatments purely by chance, with no prior sorting or grouping. Randomization alone is responsible for balancing any existing differences among units across the treatment groups.

Suppose an agricultural research station wants to compare the yield of three varieties of wheat: a standard variety currently in use, and two new varieties developed in the lab. The station has 90 available plots of farmland. In a completely randomized design, the experiment proceeds as follows:

\begin{enumerate}
    \item Measure the current condition of each plot (soil quality, drainage, sun exposure).
    \item Number the plots 1 to 90 and use a random number generator to assign 30 plots to variety 1, 30 to variety 2, and 30 to variety 3.
    \item Plant the assigned wheat variety in each plot and manage all plots identically in terms of water, fertilizer, and pest control.
    \item At harvest, measure the yield from each plot.
    \item Compare average yields across the three groups.
\end{enumerate}

% FIXME Figure/Illustration needed
% [figure placeholder: completely randomized design diagram 90 available plots randomly assigned to Group 1 (variety 1), Group 2 (variety 2), Group 3 (variety 3), all leading to Measure & Compare yield]
%\caption{Schematic diagram of a completely randomized experiment.}


This design works well when the experimental units are reasonably similar to begin with. If the 90 plots varied dramatically in soil quality, however, randomization might not perfectly balance soil quality across the three groups, particularly with a smaller experiment. This is the motivation for blocking.

\begin{Exercise}[title={Designing a Completely Randomized Experiment}, label={ex:crd}]
A sports science department wants to test whether three different warm-up routines (static stretching, dynamic stretching, or no warm-up) affect sprint times in recreational runners. Sixty volunteers have been recruited.

\begin{enumerate} 
    \item Identify the experimental units, the factor, the levels, and the treatments.
    \item Describe a complete procedure for randomly assigning the 60 volunteers to the three groups.
    \item At the end of the study, the dynamic stretching group has the fastest average sprint times. Can the researchers conclude that dynamic stretching caused the improvement? Explain.
\end{enumerate}
\end{Exercise}
\begin{Answer}[ref={ex:crd}]
\begin{enumerate} 
    \item Experimental units: the 60 volunteers. Factor: warm-up routine. Levels: three (static stretching, dynamic stretching, no warm-up). Treatments: the three levels, since there is only one factor.
    \item Number the volunteers from 01 to 60. Use \texttt{randInt(1,60,20)} on a TI-84, ignoring repeats, to select 20 volunteers for group 1 (static stretching). From the remaining 40, use \texttt{randInt} again to select 20 for group 2 (dynamic stretching). The remaining 20 form group 3 (no warm-up).
    \item Yes. Because treatments were randomly assigned, a causal conclusion is appropriate, provided the experiment was otherwise well-controlled (same track, same time of day, same measurement method, etc.). Random assignment is what licenses the claim of causation.
\end{enumerate}
\end{Answer}

\subsection{Randomized Block Design}

When experimental units differ in some important, identifiable way before the experiment begins, a completely randomized design may leave a large amount of unexplained variability in the results. A more efficient approach is to group similar units into blocks before assigning treatments.

In a \newterm{randomized block design}\index{randomized block design}, experimental units are first divided into homogeneous groups called blocks. Within each block, all treatments are randomly assigned. Because units within a block are similar to each other, the comparison between treatments within a block is cleaner than a comparison that mixes unlike units together. The blocking factor is a variable known to affect the response that the researcher controls for by separating units before randomization.

Returning to the wheat experiment: suppose the research station has two distinct soil zones, sandy soil in the northern half of the available land and clay soil in the southern half. Soil type is known to affect yield. If the three varieties are distributed randomly across all 90 plots without regard to soil type, differences in soil quality may obscure or exaggerate differences between varieties. The solution is to block on soil type.

\begin{enumerate}
    \item Divide the 90 plots into two blocks: 45 sandy-soil plots and 45 clay-soil plots.
    \item Within the sandy-soil block, randomly assign 15 plots to variety 1, 15 to variety 2, and 15 to variety 3.
    \item Within the clay-soil block, randomly assign 15 plots to variety 1, 15 to variety 2, and 15 to variety 3.
    \item Plant, manage, and harvest as before.
    \item Compare yields within each block, then examine how results differ across soil types.
\end{enumerate}

% FIXME Figure/Illustration needed
% [figure placeholder: randomized block design diagram, 90 plots separated into Sandy soil block (randomly assigned to V1, V2, V3) and Clay soil block (randomly assigned to V1, V2, V3), both leading to Measure & Compare]
% caption{Schematic diagram of a randomized block design.}

Notice that results must be compared within blocks. Comparing a variety 1 plot on sandy soil to a variety 2 plot on clay soil would confound the variety effect with the soil effect. The whole point of blocking is to make comparisons fair by ensuring the only difference within each comparison is the treatment.

A good blocking variable is one you can observe before the experiment begins and one you expect to affect the response. In the wheat experiment, soil type is visible and known to matter. In a study of a new study-skills program, prior academic performance would be a natural blocking variable. The rule: block on what you can see and control; randomize against what you cannot.

\begin{Exercise}[title={Identifying the Block Design}, label={ex:block-id}]
For each study, identify the blocking factor and explain why blocking is preferable to a completely randomized design.

\begin{enumerate} 
    \item A study tests two methods of teaching reading comprehension to primary school students. Researchers believe that students who already read above grade level will respond differently from those reading at or below grade level.
    \item An engineering team tests four surface coatings on metal panels exposed to salt water. The available panels come from two different steel mills and are known to differ in base composition.
\end{enumerate}
\end{Exercise}
\begin{Answer}[ref={ex:block-id}]
\begin{enumerate} 
    \item Blocking factor: prior reading level (above grade level vs. at or below). If above-level readers respond better to one method regardless of which method it is, mixing them randomly with below-level readers would make it harder to detect the true method effect. Blocking ensures each method is compared within similar groups.
    \item Blocking factor: steel mill source (mill A vs. mill B). Differences in base composition could affect how well coatings adhere, confounding the coating comparison. Blocking by mill source ensures all four coatings are tested on panels from each mill.
\end{enumerate}
\end{Answer}

\subsection{Matched-Pairs Design}

A \newterm{matched-pairs design}\index{matched-pairs design} is a special case of blocking in which each block contains exactly two experimental units and there are exactly two treatments to compare. 

Pairs of units are matched on some relevant characteristic, and within each pair, one unit is randomly assigned to each treatment. Because the two units in each pair are as similar as possible, any difference in their responses can be attributed to the treatment rather than to pre-existing differences.

There are two common versions of this design.

In the first, two different subjects are matched. A study comparing two physiotherapy protocols for knee rehabilitation might pair patients by age and injury severity. Within each pair, one patient is randomly assigned to protocol 1 and the other to protocol 2. The comparison is made within pairs, not across the whole sample.

In the second version, each subject serves as their own pair, receiving both treatments in a randomly determined order. This is sometimes called a \newterm{repeated measures design}\index{repeated measures design}. A washout period between treatments allows the first treatment's effect to clear before the second begins. Because each subject is compared only to themselves, all individual differences are perfectly controlled.

Consider a study testing whether caffeine affects short-term memory performance. Each participant completes a memory task twice: once after consuming a caffeinated drink, and once after consuming an identical-looking decaffeinated drink. The order is randomly determined for each participant. This is a matched-pairs design because each participant's caffeinated performance is paired with their own decaffeinated performance. Comparing one participant's caffeinated score to a different participant's decaffeinated score would not be effective.

% FIXME Figure/Illustration needed
%figure placeholder: matched-pairs diagram — participants randomly assigned order → condition 1 (caffeinated or decaf first) → washout period → condition 2 → Compare within each participant]
%{Schematic diagram of a matched-pairs design (repeated measures version).}


\begin{Exercise}[title={Matched Pairs or Randomized Block?}, label={ex:mp-block}]
Classify each study as a matched-pairs design or a randomized block design with more than two units per block. Explain your reasoning.

\begin{enumerate} 
    \item A study tests three different revision strategies (practice testing, summarizing, and re-reading) on university students. Students are divided into blocks by their first-year GPA: high (3.5 and above), mid (2.5 to 3.49), and low (below 2.5). Within each GPA block, students are randomly assigned to one of the three strategies.
    \item A study compares the effectiveness of two soil treatment methods on crop yield. Farmers are paired by farm size and climate zone. Within each pair, one farm uses method 1 and the other uses method 2.
    \item A study tests whether a standing desk reduces reported lower back discomfort. Each participant uses a standard seated desk for four weeks, then a standing desk for four weeks, with the order randomly determined for each person.
\end{enumerate}
\end{Exercise}
\begin{Answer}[ref={ex:mp-block}]
\begin{enumerate} 
    \item Randomized block design. There are three treatments and each GPA block contains more than two units. It is not a matched-pairs design because the block size exceeds two and there are more than two treatments.
    \item Matched-pairs design. There are exactly two treatments and each block (pair of farms) contains exactly two units, one receiving each treatment.
    \item Matched-pairs design (repeated measures). Each participant receives both treatments in random order and serves as their own pair of size two.
\end{enumerate}
\end{Answer}


\subsection{Choosing a Design}

Knowing what each experimental design is called is not enough. The AP exam regularly asks you to justify why a particular design is appropriate for a given situation. That requires understanding what problem each design solves and when those problems arise.

The starting point is always the completely randomized design. It is the simplest structure, and when experimental units are reasonably similar to one another, it is the right choice. Random assignment distributes any pre-existing differences among units across treatment groups by chance, and with enough units, those differences tend to balance out. If you have no strong reason to believe that a particular characteristic of the units will systematically affect the response, a completely randomized design is appropriate. 

The case for blocking arises when you can identify, before the experiment begins, a characteristic of the experimental units that you expect to affect the response variable. The problem with ignoring such a characteristic is that it introduces extra variability into your results: units in the same treatment group will differ from one another in ways that have nothing to do with the treatment, making it harder to detect a real treatment effect. Blocking solves this by grouping similar units together first, so that the comparison between treatments is made within groups of units that are alike in the relevant way. This reduces variability within treatment comparisons and makes the experiment more sensitive.

To decide whether to block, ask yourself: is there a variable I can observe right now, before any treatment is applied, that I expect to affect how units respond? If yes, and if it is practical to measure and group by that variable, block on it.

Matched pairs is the right choice when there are exactly two treatments and you can either pair units that are very similar on some relevant characteristic, or give both treatments to the same unit. Because each unit in a pair is as similar as possible to the other, the within-pair comparison is very clean. If each unit can receive both treatments (with a washout period between them), individual differences disappear entirely from the comparison.

\begin{Exercise}[title={Justifying a Design Choice}, label={ex:design-choice}]
For each situation, recommend a design (completely randomized, randomized block, or matched pairs) and write a clear justification that explains why that design is more appropriate than the alternatives.

\begin{enumerate} 
    \item Dell Technologies wants to test whether a new user interface reduces the time it takes to complete a standard task, compared to the existing interface. They have 80 volunteers available, all of whom are regular users of neither interface and have similar levels of computing experience.

    \item A sleep research lab wants to compare the effect of two room temperatures (cool vs. warm) on sleep quality. They have 40 adult volunteers. Each person's natural sleep patterns vary night to night, but researchers expect consistent individual differences in baseline sleep quality across people.

    \item Bluewave School District wants to compare three different formats for a professional development workshop (in-person, hybrid, and fully online) on teachers' reported confidence in teaching a new curriculum. The district's teaching staff includes both primary school teachers and secondary school teachers, and prior research suggests that these two groups respond very differently to professional development formats.

    \item Rhode Cosmetics wants to test whether a new moisturizer reduces skin dryness compared to a standard moisturizer. They recruit 50 volunteers and note that participants' baseline skin dryness varies widely.
\end{enumerate}
\end{Exercise}
\begin{Answer}[ref={ex:design-choice}]
\begin{enumerate} 
    \item Completely randomized design. The volunteers are similar in relevant ways (comparable computing experience, no prior exposure to either interface). There is no identified variable that would systematically affect how they respond to one interface versus the other. Randomly assign 40 volunteers to the new interface and 40 to the existing interface.

    \item Matched pairs (repeated measures). Each participant has their own baseline sleep patterns that differ from others. Using each person as their own pair eliminates individual differences in baseline sleep quality from the comparison. Randomly determine which temperature each volunteer experiences first, include a washout night between conditions, then compare each person's cool-night and warm-night sleep quality directly.

    \item Randomized block design, blocking on school level (primary vs. secondary). Prior research indicates that the two groups respond differently to professional development formats. If teachers are randomly assigned to workshop formats without blocking, differences in how primary and secondary teachers respond could mask or distort the true effect of format. Block by school level and randomly assign teachers to the three formats within each block.

    \item Randomized block design or matched pairs, blocking on baseline skin dryness level. Because baseline dryness varies widely, a participant who starts very dry might show a large reduction regardless of which moisturizer they use. Grouping participants by baseline dryness (for example, low, medium, high) and assigning both moisturizers within each group removes this variability from the comparison. Alternatively, if each participant can use each moisturizer on separate occasions, a matched-pairs repeated measures design would control individual differences even more completely.
\end{enumerate}
\end{Answer}


